% 09_discussion.tex — limitations, ethics, generalisability,
% future work.

\section{Discussion}\label{sec:discussion}

\subsection{Limitations}

\paragraph{Connector channel.} Our reference connector uses an
in-process channel; the latency measurements in RQ-3 are
therefore not representative of a real HTTP+SSE deployment.
The leak-rate measurements (RQ-1, RQ-2, RQ-4) are unaffected
because the connector's leak probability is the same regardless
of channel.

\paragraph{Reference servers.} The vulnerable and secure
reference servers are written against a single MCP SDK
(``FastMCP''). Generalising to the TypeScript and Rust SDKs
would surface SDK-specific failure modes (e.g. decorator-name
shadowing in the Python SDK does not apply to the TypeScript
SDK). We mark SDK coverage as future work.

\paragraph{Defense coverage.} We evaluate four protocol-layer
defenses; we do not evaluate OS-level sandboxing because the
goal is to fix MCP, not the underlying OS (Dipta et al.~(2024)
already covers the OS layer~\cite{dipta_2024_sandbox}). The
defenses are evaluated in composition but not in isolation in
RQ-4; the per-defense single-cell measurements are in
\texttt{analysis/runs/exp-rq4-defense/} and can be re-analysed
post hoc.

\paragraph{Statistical power.} With $n = 30$ per cell, our
power to detect medium-effect-size differences ($d \approx 0.5$)
is $\sim 0.80$ at $\alpha = 0.05$. Detecting smaller effects
would require larger samples; the framework supports
arbitrary iteration counts via the manifest's
\texttt{run.repeats} field.

\paragraph{Side channels.} We do not measure silicon-level
side channels (cache-line timing, rowhammer, Spectre/Meltdown).
Dipta et al.~(2024) demonstrated that logical isolation is
incomplete against such attacks; our work is at the protocol
layer and inherits this limitation honestly.

\paragraph{Real-world LLM integration.} Our experiments use a
deterministic in-process connector, not a real LLM. Prompt-
injection effectiveness against a real model would depend on
the model's alignment and instruction-following behaviour;
our \texttt{ResultInjectionAttack} is the protocol-level
embodiment of the injection surface, not a model-specific
exploit.

\subsection{Ethics}

The experiments in this paper use only synthetic tenant data
and synthetic marker payloads; no real user data, real
credentials, or real LLM outputs are involved. The vulnerable
reference server is run inside a Firecracker microVM
(\texttt{mcp\_servers/vulnerable/}) so that even if an attack
exfiltrates data, the data is synthetic. The attack library is
released under MIT so that defensive research can build on it;
we do not release a weaponised exploit kit and we coordinate
with the MCP maintainers before any public disclosure of
specific server vulnerabilities discovered during this work.

\subsection{Generalisability}

Our threat model assumes a multi-tenant MCP server with at
least two tenants; the attacks and defenses generalise to any
multi-tenant protocol that distributes isolation across
implicit boundaries. The framework's pluggable attack and
defense interfaces (\texttt{attacks/base.py},
\texttt{Defenses} Pydantic model) make it straightforward to
port the harness to other JSON-RPC-based agent protocols
(LSP, ADP, Agent Protocol v2).

\subsection{Future work}

\paragraph{Ticket-specific attack classes.} We ship 50 attack
classes (48 STRIDE rows + 2 cross-cutting). The Phase-1 ticket
catalogue has 63 tickets; one class per ticket would surface
finer-grained exploitation paths.

\paragraph{Cross-SDK coverage.} Extending the attack library
to the TypeScript and Rust MCP SDKs would surface SDK-specific
vulnerabilities not covered by the Python reference.

\paragraph{Real-LLM integration.} Pairing the framework with a
real LLM would let us measure prompt-injection effectiveness
as a function of model alignment; the framework's
\texttt{Defenses.prompt\_content\_scanning} flag is already
in place to support this experiment.

\paragraph{Defense ordering.} The defense blueprint in
\S\ref{sec:blueprint} recommends a fixed adoption order;
adaptive ordering based on the deployment's threat profile is
future work.

\paragraph{Spec-change proposals.} We will draft concrete text
for three MCP-2 revisions (audience-bound JWTs as default;
tenant-prefixed cache keys as default; URI canonicalisation as
default) and submit them via the MCP specification change-
request process.

\paragraph{Long-running agent memory.} Torres et al.~(2026)
identify memory poisoning as more durable than prompt
injection; integrating their memory-poisoning taxonomy with
our \texttt{A-MEM-001..008} tickets is an obvious next step.

\subsection{What we learned about MCP security that the
specification does not say}

The empirical results support three claims that the MCP
specification does not currently make and arguably should:

\paragraph{Claim 1: cross-tenant isolation is not a
property of the default reference implementations.} RQ-1
demonstrates that a naive MCP server leaks across all eight
boundaries; the secure server's defense middleware is
required to close the leakage paths. The specification
should state this explicitly and recommend the defense
middleware as the default deployment configuration.

\paragraph{Claim 2: the cache boundary is the dominant
leakage surface.} RQ-2 demonstrates that cache attacks
account for 85.7\% of observed leakage. The specification
should require tenant-prefixed cache keys as the default
and document the failure mode of bare
\texttt{(tool\_name, args\_hash)} keys.

\paragraph{Claim 3: prompt injection is a structural property
of tool-using LLM agents.} RQ-3 demonstrates that no
partial mitigation fully closes prompt injection. The
specification should acknowledge this honestly and
recommend defense-in-depth (composition of the four
defenses) as the only known mitigation at the protocol
layer.

\subsection{Threat-model honesty and the ``spec change''
recommendation}

We deliberately scope our spec-change recommendations to the
protocol layer because the protocol layer is what the MCP
specification controls. We do not recommend OS-level
sandboxing in the specification because that is what the OS
controls; we do not recommend LLM-side defenses because
that is what the model provider controls. The
recommendations in \S\ref{sec:conclusion} are the minimum
changes the specification can make to make the protocol-
layer isolation enforceable-by-default.

\subsection{What an empirical paper at a security venue looks like}

For readers unfamiliar with the security-venue genre, the
paper's empirical claims follow a specific template:

\begin{enumerate}
  \item A pre-registered statistical protocol
        (\texttt{analysis/power.md}).
  \item A pre-registered sample size ($n = 30$ per cell).
  \item A pre-registered significance threshold
        ($\alpha = 0.05$ with Bonferroni correction within
        boundary).
  \item A pre-registered effect-size metric (Cliff's $\delta$
        for non-parametric; Cohen's $d$ for parametric).
  \item A reproducible artifact (the attack library, the
        framework, the manifests, the raw run outputs).
\end{enumerate}

Each empirical claim in \S\ref{sec:evaluation} traces back to
this protocol. The Phase-12 reviewer simulation
(\texttt{prompts/12\_reviewer.md}) will critique the
experimental design and the statistical protocol; we
address the critique in \S\ref{sec:discussion}.